\documentclass[12pt]{scrartcl}
\usepackage[hagavi]{azzam}

% Induksi Matematika - Modul (gabungan Materi dan Soal Latihan)
% Materi diadaptasi dari Modul Pra-OSP Induksi (Azzam dan Akeyla):
% prinsip induksi biasa, berlangkah, dan kuat, beserta contoh jumlah kubik,
% chicken wings, dan representasi biner (ditulis ulang dengan rapi).
% Soal Latihan dari Modul Pra-OSP Induksi, esai OSP/OSN, EMC Junior, dan Klasik.
% _Solusi tetap terpisah.
% Sumber tiap soal latihan dimunculkan di PDF sebagai \item (Sumber); komentar menyimpan detail lengkap.

\title{Induksi Matematika - Modul}
\date{19 November 2026}

\begin{document}
\maketitle

\renewcommand*\contentsname{Daftar Isi}
\tableofcontents

\newpage

\part*{Materi}
\section{Induksi Biasa}

\istilah{Induksi matematika}{mathematical induction} adalah cara membuktikan
pernyataan $P(n)$ untuk setiap bilangan bulat $n$ mulai dari suatu titik awal.

\begin{theorem}[Prinsip induksi]
Misalkan $P(n)$ adalah pernyataan mengenai bilangan bulat $n$
dan $k_0$ bilangan bulat. Jika
\begin{enumerate}
    \item (\istilah{basis}{base case}) $P(k_0)$ benar, dan
    \item (\istilah{langkah induksi}{inductive step}) untuk setiap bilangan bulat $k \ge k_0$,
    $P(k)$ benar mengakibatkan $P(k+1)$ benar,
\end{enumerate}
maka $P(k)$ benar untuk setiap bilangan bulat $k \ge k_0$.
\end{theorem}

Asumsi ``$P(k)$ benar'' pada langkah induksi disebut
\istilah{hipotesis induksi}{induction hypothesis}.
Stepnya: periksa basis, asumsikan hipotesis induksi,
lalu turunkan kasus $k+1$ darinya.

\begin{example}
Buktikan bahwa
\[ 1^3 + 2^3 + \dots + n^3 = \left(\frac{n(n+1)}{2}\right)^2 \]
untuk setiap bilangan asli $n$.
\end{example}
\begin{proof}[Solusi]
Basis $n = 1$:
\[ 1^3 = 1 = \left(\frac{1 \cdot 2}{2}\right)^2. \]
Asumsikan kesamaannya berlaku untuk $n = k$. Untuk $n = k+1$,
\begin{align*}
1^3 + 2^3 + \dots + k^3 + (k+1)^3 &= \left(\frac{k(k+1)}{2}\right)^2 + (k+1)^3 \\
&= (k+1)^2\left(\frac{k^2}{4} + k + 1\right) \\
&= (k+1)^2 \cdot \frac{k^2 + 4k + 4}{4} \\
&= \left(\frac{(k+1)(k+2)}{2}\right)^2.
\end{align*}
Terbukti bahwa kesamaannya berlaku untuk $n = k+1$.
\end{proof}

\begin{example}
Buktikan bahwa $n^3 - n$ habis dibagi $6$ untuk setiap bilangan asli $n$.
\end{example}
\begin{proof}[Solusi]
Basis $n = 1$: $1^3 - 1 = 0$ habis dibagi $6$.
Asumsikan $k^3 - k$ habis dibagi $6$. Untuk $k+1$,
\begin{align*}
(k+1)^3 - (k+1) &= k^3 + 3k^2 + 3k + 1 - k - 1 \\
&= (k^3 - k) + 3k(k+1).
\end{align*}
Suku pertama habis dibagi $6$ menurut hipotesis induksi.
Hasil kali $k(k+1)$ selalu genap, jadi $3k(k+1)$ adalah kelipatan $6$.
Jumlah dua kelipatan $6$ habis dibagi $6$.
\end{proof}

\begin{example}
Argumen berikut tampak membuktikan bahwa semua kuda di dunia berwarna sama.
Temukan kesalahannya. ``Bukti.'' Induksi pada banyaknya kuda $n$.
Basis $n = 1$ jelas.
Asumsikan setiap kelompok $k$ kuda berwarna sama.
Diberikan $k+1$ kuda, keluarkan satu kuda: sisa $k$ kuda berwarna sama.
Keluarkan kuda yang berbeda: $k$ kuda sisanya juga berwarna sama.
Kedua kelompok ini saling overlap, jadi semua $k+1$ kuda berwarna sama.
\end{example}
\begin{proof}[Solusi]
Langkah induksinya gagal saat $k = 1$.
Untuk menurunkan kasus $n = 2$ dari kasus $n = 1$,
kedua kelompok satu kuda itu saling asing: tidak ada dua kuda
yang terhubung warnanya, atau tidak ada satu kuda yang mempunyai warna berbeda (yang menghubungkan keduanya).
\end{proof}

\section{Induksi Berlangkah}
\istilah{Induksi berlangkah}{induction with step} memakai ``lompatan'' $m > 1$:
kerjakan dengan $m$ basis di depan, lalu ``melompat'' sejauh $m$ setiap langkahnya.

\begin{theorem}[Prinsip induksi berlangkah]
Misalkan $P(n)$ adalah pernyataan mengenai bilangan bulat $n$,
$k_0$ bilangan bulat, dan $m \ge 1$ bilangan asli. Jika
\begin{enumerate}
    \item (Basis) $P(k_0), P(k_0 + 1), \dots, P(k_0 + m - 1)$ semuanya benar, dan
    \item (Langkah induksi) untuk setiap bilangan bulat $k \ge k_0$,
    $P(k)$ benar mengakibatkan $P(k+m)$ benar,
\end{enumerate}
maka $P(k)$ benar untuk setiap bilangan bulat $k \ge k_0$.
Prinsip ini disebut memiliki langkah $m$.
\end{theorem}

\begin{example}
Sebuah restoran menjual paket berisi $5$ chicken wing
dan paket berisi $8$ chicken wing.
Buktikan bahwa kita dapat membeli tepat $n$ chicken wing
untuk setiap bilangan asli $n > 27$.
\end{example}
\begin{proof}[Solusi]
Gunakan induksi berlangkah $5$.
Basisnya adalah lima kasus $n = 28, 29, 30, 31, 32$:
\begin{align*}
28 &= 5 + 5 + 5 + 5 + 8, \\
29 &= 8 + 8 + 8 + 5, \\
30 &= 5 + 5 + 5 + 5 + 5 + 5, \\
31 &= 8 + 8 + 5 + 5 + 5, \\
32 &= 8 + 8 + 8 + 8.
\end{align*}
Asumsikan $k$ chicken wing dapat dibeli tepat.
Untuk $k+5$, tinggal tambahkan satu paket isi $5$.
Jadi setiap $n > 27$ dapat dibeli tepat.
\end{proof}

\section{Induksi Kuat}

Pada induksi biasa, kasus $k+1$ hanya boleh memakai kasus $k$.
\istilah{Induksi kuat}{strong induction} membolehkan memakai semua kasus sebelumnya:
untuk membuktikan $P(k+1)$, seluruh $P(k_0), \dots, P(k)$ tersedia.

\begin{theorem}[Prinsip induksi kuat]
Misalkan $P(n)$ adalah pernyataan mengenai bilangan bulat $n$
dan $k_0$ bilangan bulat. Jika
\begin{enumerate}
    \item (Basis) $P(k_0)$ benar, dan
    \item (Langkah induksi) untuk setiap bilangan bulat $k \ge k_0$,
    $P(k_0), P(k_0 + 1), \dots, P(k)$ semuanya benar
    mengakibatkan $P(k+1)$ benar,
\end{enumerate}
maka $P(k)$ benar untuk setiap bilangan bulat $k \ge k_0$.
\end{theorem}

\begin{remark}
Induksi biasa adalah kasus khusus induksi kuat:
terapkan induksi kuat pada pernyataan
``$P(k_0), \dots, P(k)$ semuanya benar''.
Arah sebaliknya juga benar, jadi kedua prinsip ini setara secara logika.
Bedanya hanya kepraktisan: induksi kuat sering membuat buktinya lebih pendek.
\end{remark}

\begin{example}
Buktikan bahwa setiap bilangan asli $n$ dapat dinyatakan sebagai jumlah
bilangan-bilangan pangkat dua ($1, 2, 4, 8, \dots$) yang berbeda.
\end{example}
\begin{proof}[Solusi]
Gunakan induksi kuat. Basis $n = 1$: $1 = 2^0$.
Asumsikan setiap $1, 2, \dots, k$ mempunyai representasi demikian.
Ambil pangkat dua terbesar yang tidak melebihi $k+1$, sebut $2^l$.
Jika $k+1 = 2^l$, selesai.
Jika tidak, tulis $r = k+1 - 2^l$, sehingga
\[ 1 \le r \le k. \]
Menurut hipotesis induksi, $r = 2^{a_1} + \dots + 2^{a_m}$
dengan $a_i$ berbeda-beda. Jadi
\[ k + 1 = 2^{a_1} + \dots + 2^{a_m} + 2^l. \]
Tinggal memastikan $l$ berbeda dari semua $a_i$.
Andaikan $l = a_i$ untuk suatu $i$. Maka
\[ k+1 \ge 2^{a_i} + 2^l = 2^{l+1}, \]
padahal $2^l$ yang terbesar dengan $2^l \le k+1$, jadi $2^{l+1} > k+1$.
Kontradiksi. Akibatnya semua pangkatnya berbeda.
\end{proof}

\section{Soal Latihan}

\subsection{Induksi Biasa}
\begin{enumerate}

	% Source: Modul Pra-OSP Induksi (soal 1, bobot 1,5)
	\item 
	Buktikan bahwa
	\[ 1^2 \cdot 2 + 2^2 \cdot 3 + \dots + n^2 \cdot (n+1)
	= \frac{n(n+1)(n+2)(3n+1)}{12} \]
	untuk setiap bilangan asli $n$.

	% Source: osp smp 2003 nomor 4 esai
	\item (OSN SMP 2003)
	Buktikan bahwa $(n-1)n(n^3+1)$ senantiasa habis dibagi oleh $6$
	untuk semua bilangan asli $n$.

	% Source: Modul Pra-OSP Induksi (soal 2, bobot 1,5)
	\item 
	Buktikan bahwa
	\[ \frac{1}{2} + \frac{1}{3} + \dots + \frac{1}{2^n} \ge \frac{n}{2} \]
	untuk setiap bilangan asli $n$.

	% Source: Modul Pra-OSP Induksi (soal Identitas Catalan, bobot 2)
	\item 
	Buktikan Identitas Catalan berikut untuk setiap bilangan asli $n$:
	\[ 1 - \frac{1}{2} + \frac{1}{3} - \frac{1}{4} + \dots
	+ \frac{1}{2n-1} - \frac{1}{2n}
	= \frac{1}{n+1} + \frac{1}{n+2} + \dots + \frac{1}{2n}. \]

	% Source: osn smp 2020 nomor 6 esai
	\item (OSN SMP 2020)
	Diberikan barisan $a_1, a_2, a_3, \dots$ dengan $a_1 = 1$.
	Untuk setiap bilangan asli $2 \le n \le 2020$ berlaku
	\[ a_n = 1 + a_1 a_2 a_3 \dots a_{n-1}. \]
	Buktikan bahwa
	\[ \frac{1}{a_1} + \frac{1}{a_2} + \dots + \frac{1}{a_{2019}} < 2. \]

	% Source: Modul Pra-OSP Induksi (Kanada JMO 2020, bobot 4)
	\item (Kanada JMO 2020)
	Barisan $(a_n)$ adalah barisan bilangan real positif yang memenuhi
	$a_1 = 1$ dan $a_{n+1}^2 + a_{n+1} = a_n$ untuk setiap bilangan asli $n$.
	Buktikan bahwa untuk setiap bilangan asli $n$ berlaku
	\[ a_n \ge \frac{1}{n}. \]

	% Source: Modul Pra-OSP Induksi (soal 5, bobot 2,5)
	\item 
	Buktikan bahwa untuk setiap bilangan asli $n$,
	bilangan $3^{3^n} + 1$ merupakan hasil perkalian setidaknya $2n+1$
	bilangan prima yang tidak harus berbeda.

	% Source: Modul Pra-OSP Induksi (Malaysia MO 2018, sedikit dimodifikasi, bobot 3,5)
	% Jawaban: (a) jumlahnya sama dengan n^2; (b) 1398102
	\item (Malaysia MO 2018)
	Notasikan $g_n$ untuk $n$ sebuah bilangan asli sebagai faktor ganjil terbesar dari $n$.
	\begin{enumerate}[(a)]
		\item Buktikan bahwa $g_{n+1} + g_{n+2} + \dots + g_{2n}$
		adalah bilangan kuadrat untuk setiap bilangan asli $n$.
		\item Hitunglah nilai dari $g_1 + g_2 + \dots + g_{2048}$.
	\end{enumerate}

	% Source: Modul Pra-OSP Induksi (Klasik dan Tricky, bobot 3?)
	\item (Klasik)
	Buktikan bahwa untuk setiap bilangan asli $n$ berlaku ketaksamaan
	\[ \frac{1}{2} \cdot \frac{3}{4} \cdot \dots \cdot \frac{2n-1}{2n}
	< \frac{1}{\sqrt{3n}}. \]

	% Source: Modul Pra-OSP Induksi (soal 6, bobot 4,8)
	\item 
	Buktikan bahwa untuk setiap bilangan asli $n$,
	terdapat sebuah bilangan asli $n$ digit yang habis dibagi oleh $5^n$
	dan semua digitnya merupakan bilangan ganjil.

	% Source: emc junior 2016 (soal konstruksi bilangan rasional)
	\item (EMC Junior 2016)
	Buktikan bahwa untuk setiap bilangan asli $n$ terdapat $n$ bilangan
	rasional positif yang berbeda dengan jumlah kuadratnya sama dengan $n$.

\end{enumerate}

\subsection{Induksi Berlangkah}
\begin{enumerate}[resume]

	% Composed: induksi berlangkah 2 pada subbarisan ganjil
	% Jawaban: semua n ganjil
	\item Tentukan semua bilangan asli $n$ sehingga $2^n + 1$ habis dibagi $3$.

	% Source: Modul Pra-OSP Induksi (soal 3, bobot 2)
	\item 
	Misalkan $F_1 = 1$, $F_2 = 1$, dan $F_n = F_{n-1} + F_{n-2}$ untuk $n \ge 3$.
	Buktikan bahwa $F_{3n}$ selalu merupakan bilangan genap
	untuk sembarang bilangan asli $n$.

	% Composed: Frobenius 3 dan 5 klasik, induksi berlangkah 3
	\item Buktikan bahwa setiap bilangan asli $n \ge 8$ dapat dinyatakan
	dalam bentuk $3a + 5b$ untuk suatu bilangan cacah $a$ dan $b$.

\end{enumerate}

\subsection{Induksi Kuat}
\begin{enumerate}[resume]

	% Composed: eksistensi faktorisasi prima, induksi kuat mendasar
	\item Buktikan bahwa setiap bilangan asli $n \ge 2$ dapat dinyatakan
	sebagai perkalian bilangan-bilangan prima.

	% Source: Klasik (teorema Zeckendorf)
	\item (Klasik)
	Didefinisikan $F_1 = 1$, $F_2 = 2$, dan $F_n = F_{n-1} + F_{n-2}$
	untuk $n \ge 3$.
	Buktikan bahwa setiap bilangan asli dapat dinyatakan sebagai jumlah
	beberapa suku barisan ini yang tidak ada dua berurutan.

	% Source: osp sma 2014 nomor 5 esai
	% Jawaban: (1) Tidak prima; (2) rasionya sama dengan a_m
	\item (OSP SMA 2014)
	Misalkan $\{a_n\}$ merupakan barisan bilangan bulat yang memenuhi
	$a_1 = 2$, $a_2 = 8$, dan
	\[ a_{n+2} = 3a_{n+1} - a_n + 5(-1)^n \]
	untuk setiap bilangan asli $n$.
	\begin{enumerate}
		\item Apakah $a_{2014}$ prima?
		\item Tunjukkan bahwa untuk setiap $m$ ganjil bilangan
		\[ \frac{a_m + a_{4m}}{a_{2m} + a_{3m}} \]
		merupakan bilangan bulat.
	\end{enumerate}

\end{enumerate}

\end{document}
